% from ega4-9-fr.tex
\oldpage[IV-3]{62}

\subsection{Constructibilité de certaines propriétés des modules}
\label{subsection:IV.9.4-fr}

\begin{env}[9.4.1]
\label{IV.9.4.1-fr}
\textbf{Notations.} Dans ce numéro et les suivants jusqu'à la fin du \S~9,
nous utiliserons systématiquement les notations suivantes~: étant donné un
morphisme $f:X\to S$, nous poserons, pour tout $s\in S$,
$X_s=f^{-1}(s)=X\times_S\operatorname{Spec}(\mathbf{k}(s))$~; pour tout
$\mathcal O_X$-Module quasi-cohérent $\mathcal F$, $\mathcal F_s$ désignera
le $\mathcal O_{X_s}$-Module
$\mathcal F\otimes_{\mathcal O_S}\mathbf{k}(s)$, et pour tout homomorphisme
$u:\mathcal F\to\mathcal G$ de $\mathcal F$ dans un $\mathcal O_X$-Module
quasi-cohérent $\mathcal G$, $u_s:\mathcal F_s\to\mathcal G_s$ sera le
morphisme $p^*(u)$, en désignant par $p$ la projection canonique $X_s\to X$.
Pour toute section $g$ de $\mathcal F$ au-dessus de $X$, on désignera par
$g_s$ l'image de $g$ par l'homomorphisme canonique
$\Gamma(X,\mathcal F)\to\Gamma(X_s,\mathcal F_s)$. Pour toute partie $Z$ de
$X$, on notera $Z_s$ l'image réciproque
$p^{-1}(Z)=Z\cap X_s$ (\textbf{I}, 3.6.1). Enfin, si $Y$ est un second
$S$-préschéma et $h:X\to Y$ un $S$-morphisme, on notera $h_s$ le morphisme
$h\times1:X_s\to Y_s$.
\end{env}

\begin{proposition}[9.4.2]
\label{IV.9.4.2-fr}
Soient $S$ un préschéma intègre de point générique $\eta$, $f:X\to S$ un
morphisme de présentation finie, $\mathcal F$, $\mathcal G$, $\mathcal H$
trois $\mathcal O_X$-Modules quasi-cohérents de présentation finie. Soient
$u:\mathcal F\to\mathcal G$, $v:\mathcal G\to\mathcal H$ deux
homomorphismes de $\mathcal O_X$-Modules, et supposons que la suite
$\mathcal F_\eta\xrightarrow{u_\eta}\mathcal G_\eta
\xrightarrow{v_\eta}\mathcal H_\eta$ soit exacte. Alors il existe un
voisinage ouvert $U$ de $\eta$ dans $S$ tel que, pour tout $s\in U$, la suite
$\mathcal F_s\xrightarrow{u_s}\mathcal G_s
\xrightarrow{v_s}\mathcal H_s$ soit exacte.
\end{proposition}

Avec les notations générales de (9.2.1), il s'agit ici de la relation
$\mathbf P(X,\mathcal F,\mathcal G,\mathcal H,u,v,k)$~: «~$X$ est un
préschéma algébrique sur le corps $k$,
$\mathcal F\xrightarrow{u}\mathcal G\xrightarrow{v}\mathcal H$ une suite
exacte de $\mathcal O_X$-Modules quasi-cohérents~». Comme, pour toute extension
$k'$ de $k$, la projection canonique $X_{(k')}\to X$ est un morphisme
fidèlement plat (2.2.13), la condition (9.2.1, (i)) est satisfaite (2.2.7). En
vertu de (9.2.3), on peut se borner au cas où $S$ est intègre et
\emph{noethérien}, auquel cas $X$ est noethérien, et $\mathcal F$,
$\mathcal G$, $\mathcal H$ sont des $\mathcal O_X$-Modules cohérents.
L'hypothèse implique qu'il existe un voisinage ouvert $U$ de $\eta$ dans $S$
tel que la suite
$\mathcal F|f^{-1}(U)\to\mathcal G|f^{-1}(U)\to\mathcal H|f^{-1}(U)$ soit
\emph{exacte}, en vertu de (8.5.8, (ii)) appliqué suivant la méthode générale
de (8.1.2, \emph{a)}), et l'on peut donc déjà supposer que la suite
$\mathcal F\xrightarrow{u}\mathcal G\xrightarrow{v}\mathcal H$ est exacte~;
il suffit évidemment de prouver que l'on a
$\operatorname{Ker}(v_s)=(\operatorname{Ker}v)_s$ et
$\operatorname{Im}(u_s)=(\operatorname{Im}u)_s$ pour tout $s$ voisin de
$\eta$ dans $S$~; par suite (tenant compte de ce que les
$\mathcal O_X$-Modules $\mathcal F$, $\mathcal G$, $\mathcal H$ sont
\emph{cohérents} et de $(\mathbf 0_{\mathrm I},5.3.4)$) on est ramené à
prouver la proposition dans le cas particulier où la suite
\[
0\to\mathcal F\xrightarrow{u}\mathcal G\xrightarrow{v}\mathcal H\to0
\]
est exacte. Or, il y a alors un ouvert $U$ dans $S$ contenant $\eta$ et tel
que $\mathcal H|f^{-1}(U)$ soit \emph{plat} sur $U$ (6.9.1)~; il résulte donc
de (2.1.8) que pour tout $s\in U$, la suite
$0\to\mathcal F_s\to\mathcal G_s\to\mathcal H_s\to0$ est exacte, ce qui
achève la démonstration.

\begin{corollary}[9.4.3]
\label{IV.9.4.3-fr}
Soient $S$ un préschéma intègre, de point générique $\eta$, $f:X\to S$ un
morphisme de présentation finie. Soit
$\mathcal L^\bullet=(\mathcal L^i)_{i\in\mathbf Z}$ un complexe de
$\mathcal O_X$-Modules quasi-cohérents de présentation finie. Pour tout $i$,
il existe un voisinage ouvert $U$ de $\eta$ dans $S$ tel que les homomorphismes
canoniques
\begin{equation}
\label{IV.9.4.3.1-fr}
(\mathcal H^i(\mathcal L^\bullet))_s\to
\mathcal H^i(\mathcal L_s^\bullet) \tag{9.4.3.1}
\end{equation}
soient bijectifs pour tout $s\in U$.
\end{corollary}

